\section{The Historical Pattern of Ratcheting Emergency Authority}
\label{sec:pattern}

\subsection{The classical diagnosis: emergency as constitutive of order, or as parasitic on it}

Carl Schmitt's \emph{Political Theology}, first published in 1922 in the early years of the Weimar Republic,
opens with the assertion that ``sovereign is he who decides on the exception''
\cite{ref_schmitt_political_theology}.
The sentence is the canonical statement of a tradition that takes the state of exception to be
constitutive of sovereignty rather than parasitic on it.
On Schmitt's account, the legal order is not capable, by its own resources,
of specifying when its ordinary operation must yield to the extraordinary.
The decision is, in his terminology, a decision \emph{about} the legal order,
not a decision \emph{within} it.
The sovereign is the entity that can make this decision, and the entity whose making of it is recognized
as authoritative by the political community over which the legal order is asserted.

The Schmittian framing produces a particular relationship between the formal text of an emergency authority
and its actual exercise.
The formal text, on this account, cannot capture the exception.
The text can describe categories of circumstance, can specify procedures, can list the officials authorized to act,
and can require notification, reporting, or judicial review.
What the text cannot do is decide, in advance and from within the legal order,
whether the exception obtains.
That decision is constitutive of the order and therefore prior to it.

Giorgio Agamben's \emph{State of Exception}, published in Italian in 2003 and in English in 2005,
takes the Schmittian diagnosis and develops it into a claim about the trajectory of twentieth- and
twenty-first-century constitutional practice \cite{ref_agamben_state_of_exception}.
Agamben's argument is that the state of exception, far from being the rare and bounded condition that
classical constitutional theory contemplated, has become the dominant paradigm of governance.
The temporary suspension of ordinary law has, over a sequence of historical events,
become the permanent operating condition of the legal order.
The exception, on Agamben's reading, no longer returns the order to its ordinary shape after the
exceptional circumstance abates.
The exception persists, and the ordinary shape recedes.

The Schmitt-Agamben framing is the analytic background against which this Article writes.
The Article does not take a position on whether the political-theological reading of sovereignty is correct.
It takes Agamben's empirical claim that the exception has, in practice, ceased to return
the order to its ordinary shape,
and asks whether the persistence has a structural rather than a merely political explanation.
For Schmitt, the constitutive character of the decision is normatively defensible under conditions of genuine
emergency, not merely descriptive of how legal orders are seen to fail.
The structural argument this Article advances is therefore incompatible with that normative view,
not orthogonal to it: a Schmittian who accepts the constitutive sovereign decision will read the
typed rollback witness, and the disciplined wrapper around emergency authority more generally,
as the Kelsenian-formalist disease that the 1922 essay diagnosed, not as a corrective
\cite{ref_caldwell_popular_sovereignty,ref_dyzenhaus_legality_legitimacy}.\footnote{%
The Article is not unsympathetic to the structural reading of the political-theological tradition,
on which see also Kahn's account of political theology in liberal constitutionalism \cite{ref_kahn_political_theology}.
What follows, however, is an argument about formal grammar,
which does not require resolution of the underlying jurisprudential dispute.}

\subsection{The structural reading of the ratcheting pattern}

The historiography of the Article~48 trajectory remains contested.
Kershaw, Mommsen, and Caldwell differ on whether the 1930-32 emergency-rule period under Br\"uning
was structurally entailed by Article~48's text or politically contingent on Hindenburg's preferences
and SPD tactical toleration
\cite{ref_kershaw_hitler,ref_mommsen_weimar,ref_caldwell_popular_sovereignty}.
This Article advances the structural reading without claiming that the structural absence forces the political outcome.
The weaker claim is that the absence made the outcome structurally available, not that it made the outcome inevitable.
The converse Schmittian position canvassed in the preceding subsection,
on which decisionism is a normatively defensible response to genuine emergency rather than a pathology,
remains available to a reader who rejects the structural reading on substantive grounds.

Once the question is posed in structural terms,
the historical record yields a recurring shape that does not appear to be reducible to political contingency alone.

The shape is the following.
A grant of authority is enacted at moment $t_0$,
on the legislative understanding that the authority will be exercised during a circumstance $c$ that obtains at $t_0$,
and that the authority will lapse at moment $t_1$ when $c$ no longer obtains.
The grant of authority does not, however, include a structural mechanism that makes the lapse automatic.
The lapse is a political event, requiring affirmative action by some body
to either refrain from extending the authority or to enact its termination.
The continuation of the authority is the default;
the lapse is the deviation from the default.

The shape produces a particular kind of equilibrium failure.
At each moment $t_i > t_0$, the political question facing the body authorized to extend or terminate
the authority is the question of whether to disturb the present operating arrangement
in which the authority is active.
The disturbance carries political costs, and the maintenance of the present arrangement does not.
The expected political cost calculation favors maintenance.
Each individual maintenance decision is, in isolation, a defensible exercise of legislative discretion.
The cumulative effect of a sequence of such decisions is the transformation of an authority that was
intended as exceptional into an authority that is permanent.

Bruce Ackerman's account of the constitution of emergency, particularly in his post-September~11 writing
\cite{ref_ackerman_before_next_attack},
identifies this equilibrium failure and proposes a structural correction:
a constitutional emergency power should be enacted with a built-in supermajority requirement for extension,
such that the default at each successive period is termination rather than continuation.
Eric Posner and Adrian Vermeule, in \emph{Terror in the Balance} \cite{ref_posner_vermeule_terror},
argue that the Ackerman correction would impose costs that are not justified by the empirical record
of post-emergency governance.
Cass Sunstein's writing on sunset clauses \cite{ref_sunstein_sunset_clauses} occupies a position closer
to Ackerman than to Posner-Vermeule, holding that structural reversion to the default of termination
is a desirable feature of an emergency-power regime.

The Article takes the Ackerman-Sunstein side of this debate as background,
but adds a feature that the existing literature on sunset clauses and termination defaults does not directly address.
The feature is the requirement that the authority, at the moment of its exercise,
exhibit a constructible path back to the state of affairs that existed before the exercise.
A sunset clause requires the authority to lapse at a specified time.
It does not require that the actions taken under the authority be reversible.

The distinction is consequential.
The information collected under a surveillance authority does not return to a pre-collection state
when the authority lapses.
The content removed from a platform index under an erasure authority does not reappear in the index
when the order lapses.
The territory occupied under an authorization for use of military force is not restored to its
pre-occupation state when the authorization lapses, if it ever does.
The authorities are time-bounded in their period of grant.
They are not time-bounded in their effect.

The structural correction this Article proposes is not a sunset clause.
It is a typed rollback witness: a requirement that the authority cannot be exercised
unless, at the moment of exercise, the actor exhibits a proof that the path back to the prior state
is constructible.
The discipline applies at the moment of exercise rather than at the moment of lapse.
An authority that, at the moment of exercise, cannot exhibit a path back to the prior state
fails to type-check.
The exercise itself becomes unconstructible.

The structural correction is, in form, the same discipline that the substrate cited in
Part~\ref{sec:substrate} applies to the admission of a receipt into a polity's history.
The substrate requires, at the moment of admission, the construction of a proof that the receipt
preserves the backward-refinement property over the polity's prior history.
A receipt that does not exhibit the proof cannot be admitted.
The grammar applies at the admission step rather than at the audit step.

The conceptual move that distinguishes the typed rollback witness from a substantive prohibition
on the underlying action has a defensible lineage in Hart's analysis of legal systems
\cite{ref_hart_concept_of_law}.
Hart's central jurisprudential distinction is between primary rules, which impose duties on conduct,
and secondary rules, which are rules about rules: rules of recognition that fix the criteria of legal
validity, rules of change that govern how further rules may be enacted or amended, and rules of
adjudication that govern how the rules are to be applied.
The typed rollback witness, as this Article uses the term, is most naturally read not as a primary rule
forbidding the substance of any particular emergency action, but as a secondary rule of change:
a constraint on the form a delegated authority may take at the moment of its construction,
rather than a constraint on what the authority, once validly constructed, may substantively do.
The forbidden-versus-unconstructible distinction that this Part has drawn tracks the primary-versus-secondary
rule distinction Hart developed.
A substantively prohibited action is one that an actor must not perform; an unconstructible authority is one
that cannot validly be brought into being in the first place.
The argument operates, in Hart's vocabulary, from the internal point of view: it presupposes that legal actors
accept the structural constraint as a binding criterion governing the construction of further rules,
not merely as an external prediction about what conduct will attract sanction.
The Article does not claim that Hart's framework entails the structural-grammar argument it advances;
Hart did not address delegated emergency authority, and the application of his categories to typed witnesses
is one this Article makes rather than one it inherits.
The narrower claim is that the conceptual move from forbidden to unconstructible is not novel in jurisprudence
and stands within a tradition that the legal academy already recognizes as serious.

The remainder of the Article develops the implications of this structural correction for delegated
emergency authority,
exhibits its operation against five case studies,
and addresses the principal objections.
The objections are real and the Article acknowledges them.
The most serious objection is that the political content of the exception is not reducible to a typed
construction,
and that the formal grammar therefore misdescribes what is happening when an emergency authority is exercised.
The Article's response is that the grammar describes the wrapper, not the content,
and that even on the Schmittian view that the content is irreducibly political,
the wrapper has structure that the substrate captures.\footnote{%
The wrapper-content distinction does not resolve the underlying dispute about whether the content of the
exception is itself a matter of decision rather than law.
It allows the Article to make a narrower structural claim that is compatible with either side of that dispute.}
